\chapter{Artifacts}
The final result of the Implementation phase is a shared git-repository for the server and client.

\section{Structure}
The repository consists of three files.
\begin{itemize}
    \item \textit{server/}: Directory containing a fastapi-project for the server
    \item \textit{client/}: Directory containing an angular-project for the client
\end{itemize}

\section{Usage}
To use the generated artifacts, the following needs to be done:
\begin{enumerate}
    \item Build image for the server
    \item Build image for the client
    \item Configure the external services
    \item Start a server container
    \item Start a client container
\end{enumerate}

The following section portrays an example usage using docker-compose.

\subsection{Example Usage}
\subsection{Step 1: Building the images}
To build the images, following sequence of commands can be used on a unix-like system:
\begin{minted}[gobble=4,frame=single,linenos]{shell}
    # Building the client image
    user@host:~$ cd client
    user@host:~$ docker build . -t codex/client

    # Building the server image
    user@host:~$ cd ./../server
    user@host:~$ docker build . -t codex/server
\end{minted}

After this, the local docker system contains the images \texttt{codex/client} and \texttt{codex/client}.

\subsection{Step 2: Configuring the external services}

The following shows a \texttt{docker-compose.yml} file that declares a PostgreSQL database and a ChromaDB instance.

\begin{minted}[gobble=4,frame=single,linenos]{yaml}
    services:
        db:
            image: postgres:16-alpine
            environment:
                POSTGRES_DB: codex
                POSTGRES_USER: codex_user
                POSTGRES_PASSWORD: codex_password
            ports:
                - "5432:5432"
            volumes:
                - ./data/postgres:/var/lib/postgresql/data

        chroma:
            image: chromadb/chroma:latest
            ports:
                - "8001:8000"
            environment:
                IS_PERSISTENT: "TRUE"
                PERSIST_DIRECTORY: "/chroma/data"
                CHROMA_DB_IMPL: duckdb+parquet
            volumes:
                - ./data/chroma:/data
\end{minted}

\subsection{Step 3: Configuring the server}

The following shows the service configuration for the server in the same \texttt{docker-compose.yml} file as above.

\begin{minted}[gobble=4,frame=single,linenos]{yaml}
      server:
        image: codex/server
        depends_on:
            - db
            - chroma
        ports:
            - "8000:8000"
        volumes:
            - ./data/documents:/app/docs
            - ./data/pictures:/app/pictures
        environment:
            APP_MAIL: <mail>
            MAIL_PASSWORD: <mail-password>
            APP_MAIL_SMTP_SERVER: <smtp-host>
            APP_MAIL_SMTP_PORT: <smtp-port>
            JWT_SECRET: example-secret-here
            CLIENT_SECRET: <google-client-secret-here>
            CLIENT_ID: <google-client-id-here>
            CLIENT_BASE_URL: <url-of-client>
            DATABASE_URL: "postgresql://codex_user:codex_password@db:5432/codex"
            CHROMA_PORT: 8000
            CHROMA_HOST: chroma
            ADMIN_EMAIL: admin@example.com
            ADMIN_USERNAME: admin
            ADMIN_PASSWORD: admin
\end{minted}

\subsection{Step 4: Configuring the client}

The following shows the service configuration for the client in the same \texttt{docker-compose.yml} file as above.

\begin{minted}[gobble=4,frame=single,linenos]{yaml}
        client:
            image: codex/client
            depends_on:
                - server
            ports:
                - "8080:80"
            environment:
                SERVER_BASE_URL: <url-of-server>
                GOOGLE_CLIENT_ID: <google-client-id-here>
\end{minted}

\subsection{Step 5: Starting the containers}

To start all containers, the following bash-command is used: 

\begin{minted}[gobble=4,frame=single,linenos]{shell}
    user@host:~$ docker compose up
\end{minted}